% ======================================================================
\section{完成情况及性能参数}
% ======================================================================

\subsection{整体介绍}

见微已完成整机装配并在板上跑通完整测量流程：双路相机采集与显示、多点跟踪、停机基线、频谱判定、声学核对、阈值判级、声光与语音报警、热敏打印、JSON 记录和 APP 推送。各功能的完成情况见表~\ref{tab:status}。

\begin{table}[htbp]
\centering
\caption{功能完成情况}
\label{tab:status}
\small
\begin{tabularx}{\linewidth}{@{}L{3.0cm}L{2.0cm}Y@{}}
\toprule
\tabhead \tabh{功能} & \tabh{状态} & \tabh{依据} \\
\midrule
双路采集与 HDMI 显示 & 板上运行 & 45.0477 帧/s，452 帧无丢帧，界面实时预览 \\
多点跟踪与基线判定 & 板上运行 & 12/12 测点、11/11 窗口通过 \\
声学辅助测频 & 板上运行 & PDM 实采，声学 19.9277 Hz 与视觉一致 \\
声光与语音报警 & 板上运行 & 实测时红灯、蜂鸣器、“振动异常”语音同时动作 \\
热敏打印 & 板上运行 & 独立 UART，现场打印测量报告 \\
APP 与网关 & 板上运行 & 局域网实时数据、报警推送与结果详情 \\
双目辅助深度 & 板上运行 & PS 后端以 160×90 分辨率约每秒更新一次 \\
毫米标定 & 试验尺标 & 本次比例 0.20309 mm/像素 \\
NFC 读卡与语音播放 & 板上运行 & NFC 与 microSD 共用 SPI 总线，分别完成读卡和语音播放 \\
前面板交互与状态指示 & 板上运行 & 旋转编码器用于本机操作，5 颗 WS2812 灯显示设备状态 \\
PL 亚像素跟踪 & 核心已验证 & LK 累加核与 C++ 模型逐位一致 \\
\bottomrule
\end{tabularx}
\end{table}

\subsection{工程成果}

\subsubsection{机械成果}

V17 外壳的底壳、顶盖和 M2 垫柱三个打印件已全部打印装配，屏幕、前面板、NFC、报警灯、四个功能按键、补光灯、十字激光、相机支架和打印机均已安装到位（图~\ref{fig:mech}）。相机下倾约 25°，放在桌面时视场覆盖前方的试验台；报警灯与 NFC 位于顶盖左上，麦克风孔位于前面板右下，四个按键从左到右依次为风扇、报警灯、照明和电源允许开关。

\begin{figure}[htbp]
\centering
\begin{subfigure}[t]{0.61443\linewidth}\centering\vspace{0pt}
\includegraphics[height=74.65mm]{\figpath/photo/device_side.jpg}
\caption{整机斜视}
\end{subfigure}\hfill
\begin{subfigure}[t]{0.34557\linewidth}\centering\vspace{0pt}
\includegraphics[height=74.65mm]{\figpath/photo/interior.jpg}
\caption{壳内布置}
\end{subfigure}
\caption{整机装配成果}
\label{fig:mech}
\end{figure}

\subsubsection{电路成果}

主载板与前面板各完成一版设计、贴片和装配，装入整机运行（图~\ref{fig:circ}）。载板提供三路 5 V 与两路 3.3 V 供电，INA226 读数显示测量时整机功耗约 9.3 W；TCA9555 报警输出、DS3231 时间源以及打印机、补光、报警灯和蜂鸣器的驱动在实测中正常工作。

\begin{figure}[htbp]
\centering
\begin{subfigure}[t]{0.34562\linewidth}\centering\vspace{0pt}
\includegraphics[height=74.64mm]{\figpath/photo/carrier_wired.jpg}
\caption{主载板接线}
\end{subfigure}\hfill
\begin{subfigure}[t]{0.61438\linewidth}\centering\vspace{0pt}
\includegraphics[height=74.64mm]{\figpath/photo/touch_screen.jpg}
\caption{触摸屏与外设联调}
\end{subfigure}
\caption{电路联调}
\label{fig:circ}
\end{figure}

\subsubsection{软件成果}

PL 端完成双路视频、HDMI、PDM 接收、补光 PWM 与看门狗、打印 UART 的集成，生成可上板位流；LK 累加核、双目 Census+SGM 核、载板控制核和显示增强核完成 RTL 与仿真验证。PS 端 vibmond 完成多点测振、声学分析、双目深度、界面、记录、打印、报警、HTTP 接口和 Nano 链路，以静态 ARM 二进制部署。Nano 固件 0.2.3 已烧录，13 条语音写入 SD 卡并读回校验。APP 完成三端外壳与共用界面。

\subsection{特性成果}

\subsubsection{资源与时序}

当前位流对应的 Vivado 2025.2 实现结果见表~\ref{tab:util}。LUT 和 BRAM 约用一半，但 Slice 占用 91.3\%：视频通路的 AXI 互联、VDMA 和格式转换分散在大量 Slice 中，每个 Slice 只用了部分 LUT。Slice 的占用成为后续加入跟踪加速核时的主要布局约束。已约束路径的最差建立裕量为 +0.756 ns，位于 125 MHz 时钟域；保持裕量为 +0.028 ns。

\begin{table}[htbp]
\centering
\caption{PL 资源占用与时序（XC7Z010-1CLG400，实现后）}
\label{tab:util}
\small
\begin{tabular}{@{}lrrr@{\hspace{2.5em}}lr@{}}
\toprule
\tabhead \tabh{资源} & \tabh{使用} & \tabh{可用} & \tabh{占用率} & \tabh{时序项} & \tabh{结果} \\
\midrule
LUT & 8945 & 17600 & 50.82\% & WNS（建立） & +0.756 ns \\
触发器 & 14420 & 35200 & 40.97\% & WHS（保持） & +0.028 ns \\
Slice & 4017 & 4400 & 91.30\% & WPWS（脉宽） & +0.333 ns \\
BRAM & 30.5 & 60 & 50.83\% & 失败端点 & 0 \\
DSP & 0 & 80 & 0\% & 时钟 & 100/125/40 MHz \\
IOB & 53 & 100 & 53.00\% & MMCM & 1/2 \\
\bottomrule
\end{tabular}
\end{table}

\subsubsection{优化过程}
\label{sec:opt}

\paragraph{资源：双目深度从 PL 迁到 PS。}项目初期按“PL 做双目匹配、PS 做标定换算”的思路，在 PL 中实现了 Census + SGM 双目核并与队友原有的 BM 核做了对比。新核在 Middlebury 数据上精度明显更好\cite{scharstein2002}（表~\ref{tab:stereo}），与 C++ 模型逐位一致，板上诊断配置每 4 帧做一次 PS/PL 复算，累计 83 次核对零差异。但把双目核、DMA 和麦克风接收核一起放进整机后，Slice 接近用满，时序裕量只剩 +0.462 ns，已没有空间接入显示增强和位移核。考虑到主测振并不依赖深度，我们把双目深度降为辅助功能，移除 PL 双目核及其专用互联，改由 PS 以 160×90 分辨率约每秒计算一次。表~\ref{tab:res} 按功能阶段列出了资源变化：LUT 减少 5081 个，BRAM 减少 11 块，DSP 从 3 个降到 0。表中同时列出各阶段保留的功能。这里统计的是整机实现的净变化，包括双目核、专用 DMA、AXI 互联和重新布局产生的变化。表~\ref{tab:stereo} 中的核级综合估计对应单独的计算模块。

\begin{table}[htbp]
\centering
\caption{PL 资源随架构调整的变化（整机实现统计）}
\label{tab:res}
\small
\begin{tabularx}{\linewidth}{@{}L{3.2cm}Yrrrr@{}}
\toprule
\tabhead \tabh{阶段} & \tabh{功能范围} & \tabh{LUT} & \tabh{BRAM} & \tabh{DSP} & \tabh{Slice} \\
\midrule
PL 双目阶段 & 双路视频、PL 双目核 + DMA、PDM 接收 & 13655 & 41.5 & 3 & 4400 \\
PS 深度阶段 & 移除 PL 双目核，深度迁到 PS & 8574 & 30.5 & 0 & 4026 \\
当前位流 & 在上一阶段基础上增加打印 UART & 8945 & 30.5 & 0 & 4017 \\
\bottomrule
\end{tabularx}
\end{table}

\begin{figure}[htbp]
\centering
\begin{tikzpicture}
\begin{axis}[width=0.86\linewidth,height=5.6cm,ybar,bar width=13pt,ymin=0,ymax=112,
  symbolic x coords={PL 双目阶段,PS 深度阶段,当前位流},xtick=data,enlarge x limits=0.22,
  ylabel={占用率 / \%},legend style={at={(0.5,1.03)},anchor=south,legend columns=3,font=\footnotesize,draw=none},
  tick label style={font=\footnotesize},label style={font=\footnotesize},nodes near coords,
  nodes near coords style={font=\tiny},every node near coord/.append style={/pgf/number format/.cd,fixed,precision=1},
  grid=major,grid style={draw=gray!20}]
\addplot[fill=axred!70,draw=axred!60!black] coordinates {(PL 双目阶段,77.6) (PS 深度阶段,48.7) (当前位流,50.8)};
\addplot[fill=navy!55,draw=navy] coordinates {(PL 双目阶段,69.2) (PS 深度阶段,50.8) (当前位流,50.8)};
\addplot[fill=amber!80,draw=amber!60!black] coordinates {(PL 双目阶段,100.0) (PS 深度阶段,91.5) (当前位流,91.3)};
\legend{LUT,BRAM,Slice}
\end{axis}
\end{tikzpicture}
\caption{三个阶段的 PL 资源占用率}
\label{fig:res}
\end{figure}

\begin{table}[htbp]
\centering
\caption{PL 双目核与原 BM 核的对比（9 个 Middlebury 场景，模拟两相机曝光差异，$D=64$）}
\label{tab:stereo}
\small
\begin{tabular}{@{}lccc@{}}
\toprule
\tabhead \tabh{方案} & \tabh{有效像素密度} & \tabh{误差 >2 像素比例} & \tabh{RMS 误差} \\
\midrule
原 BM 核 + 空间支持滤波 & 77.3\% & 8.2\% & 5.42 像素 \\
新 PL 核原始输出 & 90.9\% & 4.4\% & 2.10 像素 \\
新 PL 核 + PS 后处理 & 87.0\% & 3.1\% & 1.46 像素 \\
\midrule
\multicolumn{4}{@{}p{0.86\linewidth}@{}}{资源（综合估计）：新核约 2.0k LUT、1.5k FF、5.5 BRAM、4 DSP；原 BM 核 6605 LUT、4442 FF、17 BRAM} \\
\bottomrule
\end{tabular}
\end{table}

\paragraph{精度：跟踪算法从 NCC 改为逆组合 LK。}早期单点跟踪使用整数 NCC 搜索加抛物线插值。在 32×32 模板、斜坡加正弦的合成亚像素平移序列上，NCC 抛物线插值的误差 RMS 为 0.13 像素，存在明显的系统偏差（像素锁定）；逆组合 LK 在无噪声时为 0.005 像素，加 $\sigma=2$ 灰度噪声时为 0.014 像素，精度提高 10～25 倍。多点测振核因此采用 LK，NCC 仅保留为质量判据。

\paragraph{实时性：面向 PL 的 LK 累加核与 PS 端的等价优化。}LK 每次迭代的主要计算是双线性插值和五个累加和（$\sum I$、$\sum I^2$、$\sum g_xI$、$\sum g_yI$、$\sum TI$）。我们在 PL 中实现了 \code{lk\_accum}：每个时钟处理一个像素，插值精度 1/256 像素，五路 48 位整数乘加，整数 ZNCC 搜索与 LK 亚像素更新共用同一个引擎。与 C++ 位精确模型比较，合成纹理 8 例加实拍帧 8 例共 16 组黄金向量逐位一致，32×32 模板每遍 1029 个时钟。按每测点每帧约 81 次整数搜索、最多 25 次细搜和 6 次 LK 迭代估算，100 MHz 下单测点约 1.2 ms，4 个测点加 1 个背景约 5.8 ms，低于 45 帧/s 的 22 ms 帧周期，为后续把多点跟踪迁入 PL 准备好了核心。PS 端在不改变数值结果的前提下做了两项优化：高斯平滑只处理模板可能读到的区域（约为整帧计算量的十分之一）；模板插值时预先缓存 25 个行坐标与 25 个列坐标的整数部分和小数部分，插值与累加顺序保持不变，结果与原核逐位一致。

\paragraph{鲁棒性：采样、参考点与打印流程。}频率搜索采用 30 Hz 软件上限，并结合实测帧率与奈奎斯特边界检查，7、21.6、22 Hz 合成信号回归通过。短缺帧按连续不超过 2 帧、总量不超过 5\%、间隔偏差不超过 10\% 的条件重采样，记录保留原始帧数与丢帧标志。静止参考点选在清晰二维纹理区域，69 个基线窗口全部通过。打印采用 24 行分段、UART 排空和 120 ms 间隔，同一份 720 行报告正常输出。表~\ref{tab:optsum} 汇总优化前后的处理策略与结果。

\begin{table}[htbp]
\centering
\caption{优化前后对比}
\label{tab:optsum}
\small
\begin{tabularx}{\linewidth}{@{}L{2.8cm}L{3.6cm}L{3.6cm}Y@{}}
\toprule
\tabhead \tabh{项目} & \tabh{优化前} & \tabh{优化后} & \tabh{依据} \\
\midrule
PL LUT / BRAM / DSP & 13655 / 41.5 / 3 & 8945 / 30.5 / 0 & 实现报告，功能范围不同 \\
跟踪误差（合成） & NCC 0.13 像素 & LK 0.005～0.014 像素 & 32×32 合成序列 \\
双目坏点率 & BM 8.2\% & 新核 + 后处理 3.1\% & Middlebury 9 场景 \\
可测频率上限 & 20.3 Hz（$0.45f_s$） & 22.5 Hz（奈奎斯特） & 22 Hz 合成回归 \\
丢帧处理 & 任一丢帧整窗无效 & ≤5\% 有界重采样 & 回归测试与记录标志 \\
平滑计算量 & 整帧 & 约 1/10，结果逐位一致 & 模板读取范围分析 \\
PDM 抽取 & 逐位积分（约 1000 万次加法/s） & 查表（125 万次/s） & 与逐位实现逐位一致 \\
打印 & 连续光栅传输 & 24 行分段，720 行报告正常 & 同一报告重打对照 \\
\bottomrule
\end{tabularx}
\end{table}

\subsubsection{测试与验证}

各部分的测试方式与结果见表~\ref{tab:tests}。所有比较类测试都以退出码和逐项差异判定，不以“程序没有报错”作为通过依据。

\begin{table}[htbp]
\centering
\caption{测试与验证汇总}
\label{tab:tests}
\small
\begin{tabularx}{\linewidth}{@{}L{3.6cm}YL{3.4cm}@{}}
\toprule
\tabhead \tabh{对象} & \tabh{方法} & \tabh{结果} \\
\midrule
PDM 接收核 & Icarus 测试平台，寄存器读写与 FIFO 行为 & 通过 \\
PDM 抽取 & 查表实现与逐位 CIC/FIR 参考对比，随机与边界码流 & 逐位一致 \\
PDM 寄存器表 & 单元测试解析 RTL 读写译码与驱动头文件比对 & 17 个寄存器一致 \\
LK 累加核 & 16 组黄金向量（合成 8 + 实拍 8） & 16/16 逐位一致 \\
双目核（PL 深度阶段） & RTL/C++ 逐位对比；板上每 4 帧复算 & 83 次零差异 \\
测振算法 & 合成位移回归：7/21.6/22 Hz、两种相位、丢帧与采样率失配 & 通过，失配被拒绝 \\
报警逻辑 & 主机测试：判级、静音、重试、刷新、阈值修改重判、无效结果清灯 & 通过 \\
主机集成 & PS 主程序与 ESP32 网关核心集成检查 & 105 项通过 \\
ARM 构建 & 静态链接 vibmond；代码包保存部署快照与 SHA256 & 已归档 \\
Nano 固件 & 0.2.3 编译、DFU 烧录、USB 命令应答、SD 语音读回校验 & 通过 \\
板端实测 & 2026-11-04 完整测量 & 见表~\ref{tab:record} \\
技能包 & 一键离线复跑：正例、反例、故障注入 & 见 \ref{sec:llm} 节 \\
\bottomrule
\end{tabularx}
\end{table}

\subsection{复现说明}

完整工程发布在 \url{https://github.com/quanxian318/jianwei-vibrometer}。仓库按源码、仿真、构建、板端文件、测试数据、技能包和报告组织，同时提供硬件工程、外壳模型、演示视频及中英文海报。主要入口见表~\ref{tab:repo}。

\begin{table}[htbp]
\centering
\caption{代码与工程入口}
\label{tab:repo}
\small
\begin{tabularx}{\linewidth}{@{}L{6.4cm}Y@{}}
\toprule
\tabhead \tabh{目录或文件} & \tabh{内容} \\
\midrule
\path{src/ax7010/linux/vibmon} & PS 主程序、测振、声学、界面、工具与测试 \\
\path{src/ax7010/pl} & RTL、测试平台与黄金向量 \\
\path{build/vivado} & 当前 BD、RTL、用户 IP、约束和实现报告 \\
\path{build/vivado_tests} & 五项分项工程、前面板测试源码、位流和报告 \\
\path{src/nano_esp32} & Nano 固件、预编译文件与 SD 卡语音 \\
\path{src/mobile_app}、\path{src/shared} & 多端 APP 与共用协议 \\
\path{sim}、\path{data} & 仿真入口、测试结果与板端测量 JSON \\
\path{board/artifacts} & PL 位流与 PS 可执行快照 \\
\path{hardware} & 原生 EDA 工程、原理图、BOM、网表及外壳模型 \\
\path{docs} & 接线、模块设置、构建说明及散斑图样 \\
\path{skill}、\path{report} & 技能包、协作记录和技术报告 \\
\path{media} & 演示视频、海报和产品图片 \\
\bottomrule
\end{tabularx}
\end{table}

\paragraph{1. 硬件装配。}用嘉立创 EDA 专业版打开原生工程，按 Gerber 与 BOM 制作两块 PCB；外壳采用最终模型中的底壳、顶盖和垫柱三个打印件。按照逐引脚接线表连接 AX7010、主载板、前面板及外设。PN532 使用 SPI 模式，拨码开关 1 置 OFF、2 置 ON；MP34DT01 的 L/R 接地。接入 AX7010 前检查电源输出和线束针序。

\paragraph{2. PL 构建。}在仓库根目录打开终端，使用 Vivado 2025.2：
\begin{lstlisting}
cd build/vivado
vivado -mode batch -source create_project.tcl -tclargs build
\end{lstlisting}
脚本指定器件 \code{xc7z010clg400-1}，加载 BD、RTL、约束与用户 IP 后执行综合、实现和位流生成。归档的实现结果保存在 \path{reports_original_build}。

\paragraph{3. PS 构建与测试。}在 Linux 或 WSL 中，从仓库根目录执行：
\begin{lstlisting}
cd src/ax7010/linux/vibmon
make -j4 BUILD=build_host \
  build_host/test_drivers build_host/test_vibration_band \
  build_host/test_multipoint build_host/test_render \
  build_host/test_uart16550
build_host/test_drivers
build_host/test_vibration_band
build_host/test_multipoint
build_host/test_render
build_host/test_uart16550
\end{lstlisting}
交叉编译 ARM 程序：
\begin{lstlisting}
make -j4 BUILD=build_arm \
  CXX=arm-linux-gnueabihf-g++ CC=arm-linux-gnueabihf-gcc \
  AR=arm-linux-gnueabihf-ar \
  CXXFLAGS='-O2 -g -mfpu=neon -mfloat-abi=hard' \
  CFLAGS='-O2 -g -mfpu=neon -mfloat-abi=hard' \
  LDFLAGS='-static -pthread' build_arm/vibmond
\end{lstlisting}

\paragraph{4. 板端运行。}系统使用 ALINX AX7010 的 512 MB DDR Linux 镜像。视频缓存占用 DDR 顶端 32 MiB，预留范围为 \code{0x1E000000--0x1FFFFFFF}。先在主机上处理基础镜像中的设备树，将输出文件作为启动分区的 \code{image.ub}：
\begin{lstlisting}
python3 src/ax7010/linux/vibmon/tools/prepare_an5642_image.py \
  image_original.ub image.ub
\end{lstlisting}
板卡启动后停止相机和界面进程，加载 PL：
\begin{lstlisting}
python3 src/ax7010/linux/vibmon/tools/prepare_an5642_runtime.py \
  board/artifacts/printer_uart.bit.bin
\end{lstlisting}
该脚本配置 FCLK1 为 125 MHz、FCLK2 为 40 MHz，并开启两路 PS I²C 时钟。Nano 使用 PS UART0 的 EMIO 接口，示例设备节点为 \code{/dev/ttyPS1}，设置为 115200、8N1、RTS/CTS 硬件流控。板级配置以 \path{config/current_board.example.json} 为入口，将相机、显示、串口、I²C 和标定文件路径设置为板上的实际路径后保存为 \code{board.json}。

源码构建程序使用 \code{layout\_hash} 绑定毫米标定与测点布局。标定文件记录尺标比例、方向和布局哈希，三者匹配后输出毫米值。归档程序快照 \path{board/artifacts/vibmond.gz} 保留当次实测部署行为，沿用当前尺标比例。复现实测快照时，解压并运行：
\begin{lstlisting}
gzip -dk board/artifacts/vibmond.gz
chmod +x board/artifacts/vibmond
board/artifacts/vibmond -c board.json
\end{lstlisting}
重新编译源码后采用布局绑定标定；需要复现快照中的比例保留方式时，使用 \path{tools/reproduce_legacy_binary.py}，命令见仓库板端说明。

\paragraph{5. Nano 与 APP。}从仓库根目录使用 arduino-cli 编译并烧录 Nano 固件：
\begin{lstlisting}
arduino-cli compile --fqbn arduino:esp32:nano_nora \
  src/nano_esp32/vibmon_nano
arduino-cli upload --fqbn arduino:esp32:nano_nora \
  -p <port> src/nano_esp32/vibmon_nano
\end{lstlisting}
将 \path{src/nano_esp32/sd_card/voice} 下的 13 条 WAV 复制到 FAT32 microSD 卡的 \path{/voice} 目录。APP 共用网页位于 \path{src/mobile_app/web}，Android、iOS 和 Windows 外壳分别使用 Android Studio、Xcode 和 Visual Studio 构建。

\paragraph{6. 测量。}在被测面贴散斑，参考点选在静止背景的清晰纹理处。固定机位后设置测量框和参考点，完成尺标标定，停机采集约 15 s 基线，启动被测物后测量约 10 s。结果页显示主频、位移峰峰值、有效值、声学核对和报警状态，可在 APP 中查看并打印报告。更换机位后重新标定、采集基线。

\FloatBarrier

